\begin{titlepage}
\begin{tikzpicture}[overlay, remember picture, shift={(current page.center)}, yshift=0.1\paperheight]
    \def\Ox{0}
    \def\Oy{0}
    
    \draw[blue!30, line width=0.3pt] (\Ox,\Oy) circle (8cm);
    \draw[blue!20, line width=0.3pt] (\Ox,\Oy) circle (6.5cm);
    \draw[blue!10, line width=0.3pt] (\Ox,\Oy) circle (5cm);
    
    \def\R{8}
    \def\omega{30}
    \pgfmathsetmacro{\theta}{\omega}
    \pgfmathsetmacro{\Px}{\R*cos(\theta)}
    \pgfmathsetmacro{\Py}{\R*sin(\theta)}
    
    \filldraw[red!60] (\Px,\Py) circle (2.5pt);
    
    \def\r{1.8}
    \def\omegai{120}
    \pgfmathsetmacro{\thetai}{\omegai}
    \pgfmathsetmacro{\Pix}{\Px + \r*cos(\thetai)}
    \pgfmathsetmacro{\Piy}{\Py + \r*sin(\thetai)}
    
    \draw[red!50, line width=0.4pt] (\Px,\Py) circle (\r);
    \filldraw[red!50] (\Pix,\Piy) circle (2.0pt);
    
    \def\rr{0.8}
    \def\omegaii{210}
    \pgfmathsetmacro{\thetaii}{\omegaii}
    \pgfmathsetmacro{\Piix}{\Pix + \rr*cos(\thetaii)}
    \pgfmathsetmacro{\Piiy}{\Piy + \rr*sin(\thetaii)}
    
    \draw[orange!70, line width=0.4pt] (\Pix,\Piy) circle (\rr);
    \filldraw[orange!80] (\Piix,\Piiy) circle (1.6pt);
    
    \def\rrr{0.35}
    \def\omegaiii{330}
    \pgfmathsetmacro{\thetaiii}{\omegaiii}
    \pgfmathsetmacro{\Piiix}{\Piix + \rrr*cos(\thetaiii)}
    \pgfmathsetmacro{\Piiiy}{\Piiy + \rrr*sin(\thetaiii)}
    
    \draw[green!60!black, line width=0.4pt] (\Piix,\Piiy) circle (\rrr);
    \filldraw[green!70!black] (\Piiix,\Piiiy) circle (1.2pt);
    
    \draw[blue!15, line width=0.2pt, densely dotted] (\Ox,\Oy) -- (\Px,\Py);
    \draw[red!15, line width=0.2pt, densely dotted] (\Px,\Py) -- (\Pix,\Piy);
    \draw[orange!15, line width=0.2pt, densely dotted] (\Pix,\Piy) -- (\Piix,\Piiy);
    \draw[green!15, line width=0.2pt, densely dotted] (\Piix,\Piiy) -- (\Piiix,\Piiiy);
    
    \foreach \k in {0,1,...,59}
    {
        \filldraw[blue!40] ({360/60*\k}:8cm) circle (1.2pt);
        \filldraw[blue!30] ({360/60*\k + 3}:6.5cm) circle (1.0pt);
        \filldraw[blue!20] ({360/60*\k + 6}:5cm) circle (0.8pt);
    }

    \foreach \k in {0,1,...,11}
    {
        \filldraw[red!50] ({360/12*\k}:9cm) circle (2.0pt);
    }

    \foreach \k in {0,1,...,7}
    {
        \filldraw[green!50!black] ({360/8*\k + 11.25}:7.2cm) circle (1.8pt);
    }

    \foreach \k in {0,1,...,4}
    {
        \filldraw[orange!60] ({360/5*\k + 18}:5.8cm) circle (1.6pt);
    }
\end{tikzpicture}

\begin{tikzpicture}[overlay, remember picture, shift={(current page.south)}, yshift=0.5cm]
    \node[font=\footnotesize, text=black!30] at (0,0) 
        [yshift=2.0cm] {$\mathcal{P}_T = \{(A_k, \omega_k, \phi_k) \mid A_k \in \mathbb{R}^+, \omega_k = 2\pi k/T, \phi_k \in [0, 2\pi), k = 0, 1, \dots, T-1\}$};
    
    \node[font=\footnotesize, text=black!30] at (0,0) 
        [yshift=1.5cm] {$\mathcal{K}^{(d)}\{f, g\} = \sum_{n}\sum_{m} A_{n,m} \cdot \operatorname{Similarity}^{(d)}(\theta_n, \theta_m) \cdot e^{i(\Omega_{n,m} x + \Phi_{n,m})}$};
\end{tikzpicture}

\vspace*{0.5cm}

\begin{center}

{\fontsize{36pt}{42pt}\selectfont \textbf{周期数学}}

\vspace{0.1cm}

\begin{center}
\itshape
“道可道也，非常道也。名可名也，非常名也。” \\[0.2cm]
“道生一，一生二，二生三，三生万物。” \\[0.2cm]
——《道德经》
\end{center}

\vspace{0.1cm}

\begin{center}
\itshape
“观天之道，执天之行，尽矣。” \\[0.2cm]
——《阴符经》
\end{center}

\vspace{0.1cm}

\begin{center}
\itshape
“数学是上帝用来书写宇宙的语言。” \\[0.2cm]
——伽利略·伽利莱
\end{center}

\vspace{0.1cm}

\begin{center}
\itshape
“这个发现让我确信，伏羲不仅是中国哲学的创始人，而且是全人类最古老的智者之一。他所留下的符号系统，包含了一种普遍的文字——一种能够表达所有思想的字符。” \\[0.2cm]
“中国人似乎在远古时代就已经拥有了一种普遍的文字，八卦就是它的遗迹。只是后来的人忘记了它的真正含义，只把它用于占卜。” \\[0.2cm]
——莱布尼茨
\end{center}

\vspace{0.1cm}

\begin{center}
\itshape
“同声相应，同气相求。” \\[0.2cm]
“无平不陂，无往不复。” \\[0.2cm]
——《易经》
\end{center}

\vspace{0.1cm}

\begin{center}
\itshape
“虚霩生宇宙，宇宙生气，气有涯垠，清阳为天，重浊为地，天地袭精为阴阳，阴阳专精为四时，四时散精为万物。” \\[0.2cm]
——《淮南子·天文训》
\end{center}

\vspace{0.1cm}

\begin{center}
\itshape
"方生方死，方死方生；方可方不可，方不可方可。因是因非，因非因是。" \\[0.2cm]
——《庄子·齐物论》
\end{center}

\begin{center}
\itshape
“大同而与小同异，此之谓小同异；万物毕同毕异，此之谓大同异。” \\[0.2cm]
——《庄子·天下》
\end{center}

\vfill
\end{center}
\end{titlepage}

\cleardoublepage

